\renewcommand{\stagefigureprefix}{2}\setcounter{figure}{0}
\reporttitle{第二阶段：指令设计与基础测试}

\section{设计目标与实施范围}
点积把两组数据逐项相乘，再将乘积求和，是量化矩阵运算的重要计算步骤。本设计在昆明湖V2已有向量路径中增加八项INT8点积指令，将乘法与归约组合为专用数据通路，同时复用处理器的译码、寄存器重命名、调度、写回和精确提交机制。

通用构建默认采用单请求实现，弹性两槽实现作为可选优化保留。本章重点介绍两槽实现的算术结构、背压和取消处理，并通过正数、负数及边界值测试展示结果。第三阶段的主要周期和成本数据对应明确标识的八项两槽受测版本；十六项扩展的最新功能验证列在本章末尾。这样的组织覆盖赛题2.1的设计说明、2.2的RTL实现和2.3的三类基础测试。

\section{指令语义与编码}
\subsection{八项有符号点积与结果布局}
本指令采用128位向量寄存器（VLEN=128）、8位元素（SEW=8）和单寄存器分组（LMUL=1）；有效元素数为8（VL=8），起始执行索引为0（vstart=0）。指令读取两个源寄存器的低8个字节，将每个字节解释为有符号整数，计算
\[
S=\sum_{i=0}^{7}\mathrm{int8}(vs1[8i+7:8i])\,
                    \mathrm{int8}(vs2[8i+7:8i]).
\]
目的低32位保存$S$的二补码，高96位写零。运算从零开始累加，读取指令执行前的源值，因此支持目的与源寄存器重合，也支持使用v0作为普通源或目的。每个源以128位向量传入数据通路，八项算术使用其中低64位的8个元素。

\reportfigure{S2-01}{八项点积指令字段与完整结果布局}{11.0cm}{两个源各提供8个有符号字节；目的低32位保存点积结果，高96位写零。}{设计示意，依据指令规格与实际译码字段。}

\subsection{数值范围与中间位宽}
有符号8位数的范围为$-128$至127，单个乘积范围为$[-16256,16384]$，八项和范围为$[-130048,131072]$。乘积使用16位；按二项、四项、八项归并时，精确位宽依次增加为17、18、19位。正端点131072比18位有符号数的最大值大1，因此极值用例能直接检验归约过程中是否保留了完整进位。

归约结果先符号扩展为32位，再与96位零拼接。负数的符号由低32位字段表达，例如$-8$的低32位为\texttt{0xfffffff8}，完整128位结果仍以96个零作为高位。内部精确和的19位、结果槽的32位和公开写回的128位分别承担算术、存储和架构接口三种作用。

\subsection{编码与合法执行状态}
指令采用custom-0编码，\texttt{MATCH=0x0200000b}、\texttt{MASK=0xfe00707f}。
\begin{center}
\small
\begin{tabular}{lll}
\toprule
字段 & 位段 & 含义或取值\\
\midrule
funct6 & 31:26 & \texttt{000000}\\
vm & 25 & \texttt{1}\\
vs2、vs1 & 24:20、19:15 & 两个向量源寄存器编号\\
funct3 & 14:12 & \texttt{000}\\
vd & 11:7 & 向量目的寄存器编号\\
opcode & 6:0 & custom-0，\texttt{0x0b}\\
\bottomrule
\end{tabular}
\end{center}
汇编记法为\texttt{vdot.vv vd, vs2, vs1}。

例如\texttt{vdot.vv v10,v8,v9}对应\texttt{0x0284850b}。程序通过经过校验的编码宏或机器字生成该指令。

向量功能启用、vill=0、vm=1以及本节的SEW、LMUL、VL和vstart共同确定合法状态，四种vta/vma策略组合均适用。合法执行完成一次完整目的写回，并保持VL、vtype、vlenb、vxrm、vxsat和vcsr。非法请求携带illegal-instruction异常完成，同时关闭向量与v0目的写使能，由处理器的精确异常机制处理。

\section{数据通路与弹性两槽控制}
\subsection{乘积槽与结果槽}
点积数据通路包含8个有符号$8\times8$乘法器、乘积槽P、归约网络和结果槽R。P保存8个16位乘积及这笔请求的控制、性能信息；R保存32位结果及同一请求的身份。结果从R输出时补入96位零，兼顾紧凑的内部状态和完整的向量写回接口。

\reportfigure{S2-02}{八项点积的乘积槽、归约网络和结果槽}{11.0cm}{P保存8路16位乘积，R保存32位结果；控制与目的身份随数据流动，输出接口为128位。}{设计示意，依据八项弹性FU实现及32位结果保存结构。}

两槽的有效位表示各自是否保存一笔请求。重排序缓冲区（ROB）编号、目的物理寄存器、向量或v0写使能、异常信息以及性能字段与数据同时移动，使写回始终能够找到正确的目的。输出停顿时，新输入的变化与R中已有请求相互独立。

\subsection{时钟、复位与连续流动}
VDot功能单元（FU）接入共享执行单元VFEX4，后文将VFEX4称为父级。实现沿用功能单元所在的单时钟域和父级复位连接。P、R在时钟上升沿更新；复位清除两槽有效位，并抑制复位期间的输入和输出握手。在父级选择并允许接收VDot、独立请求连续供给、下游持续就绪且没有复位或重定向时，请求于E0进入P，E1完成归约并进入R，E2是最早的输出握手沿。连续请求的发起间隔为1拍。

\begin{center}
\small
\begin{tabular}{p{0.16\linewidth}p{0.16\linewidth}p{0.16\linewidth}p{0.40\linewidth}}
\toprule
P状态 & R状态 & 下游状态 & 本拍动作\\
\midrule
空 & 空 & 任意 & 接受新请求到P\\
有效 & 空 & 任意 & 原P进入R，同时可接收下一请求\\
空 & 有效 & 反压 & 保持R，P可接收一笔请求\\
有效 & 有效 & 反压 & 保持两槽，暂停接收\\
有效 & 有效 & 就绪 & 原R输出、原P进入R、新请求进入P\\
\bottomrule
\end{tabular}
\end{center}

\reportfigure{S2-05}{连续请求在两槽中的逐拍流动}{8.0cm}{持续供给且下游就绪时，旧R输出、旧P推进和新请求接收可以在同一拍重叠；最早响应为接受后2拍，发起间隔为1拍。}{逐拍机制示意，依据弹性槽状态方程和连续请求验证。}

子单元就绪与父级准入分别判定。槽位腾出提供接收容量，父级仍根据当前选择与共享资源占用决定何时派发新请求。

时延描述一笔请求从接受到响应的距离，发起间隔描述相邻请求的接受节奏。两槽结构支持连续独立点积；程序的最终周期还由数据依赖、调度、访存和写回共同决定。功能单元继续通过valid/ready报告完成，反压时可自然延后响应。

\subsection{背压、重定向与事务守恒}
R有效而下游未就绪时，数据、目的和控制包一起保持。P可利用空槽再接收一笔请求；若两槽均占用且R仍受阻，输入ready撤销，直到R输出或取消腾出空间。

\reportfigure{S2-06}{背压期间的数据与目的身份保持}{8.0cm}{R在输出阻塞期间保存完整响应；空闲P可继续接收一笔请求，满槽后向上游施加反压。}{控制机制示意，依据FU和真实父级反压验证。}

重定向按两个槽各自的ROB年龄独立判断取消，保留较老的幸存请求，并在该拍暂停新请求接收。被取消的R停止输出，被取消的P释放待归约状态；复位的优先级最高。ROB环回场景使用完整指针年龄关系确定先后。

\reportfigure{S2-07}{部分取消、全部取消与复位恢复}{8.0cm}{分别判断P、R中的请求是否被重定向命中，展示幸存请求、取消清空和复位后重新接收。}{控制机制示意，依据槽内ROB判断和对应定向用例。}

每次验证在统一时钟坐标下检查
\[
N_{\mathrm{accept}}=N_{\mathrm{response}}+N_{\mathrm{cancel}}
                  +N_{\mathrm{reset}}+N_{\mathrm{inflight}},
\qquad 0\leq N_{\mathrm{inflight}}\leq2.
\]
该关系将每笔已接受请求归入响应、取消、复位丢弃或当前占用之一，并在运行结束时检查两槽排空。每笔响应还核对完整128位结果、ROB和目的身份。

\reportfigure{S2-10}{两槽同拍推进与取消、复位处理}{10.0cm}{满槽时原R输出、原P推进与新请求接受可在同一拍进行；取消与复位示例说明被撤销请求和幸存请求的处理。}{状态机制示意，依据实际有效位更新及事务守恒检查。}

\section{RTL实现与处理器接入}
\subsection{关键实现语句}
算术部分位于\seqsplit{src/main/scala/xiangshan/backend/fu/vector/VDotDatapath.scala}，乘积和每级归约按有符号类型计算，并逐级保留进位。\seqsplit{wrapper/VDotElastic.scala}保存P、R的控制和数据。下面摘示从P转入R的关键语句：
\begin{Verbatim}[fontsize=\small,breaklines=true]
when(move) {
  rPacket.ctrl := pCtrl
  rPacket.perfDebugInfo := pPerf
  rPacket.res.data := Mux(pCtrl.exceptionVec.get(ExceptionNO.illegalInstr),
    0.U(32.W), reduce.io.result(31, 0))
}
\end{Verbatim}
这些赋值读取本拍边沿前的P内容，因此同拍接收的新请求会进入P，原请求的结果和身份则进入R。公开输出使用：
\begin{Verbatim}[fontsize=\small,breaklines=true]
io.out.bits.res.data := Cat(0.U(96.W), rPacket.res.data)
\end{Verbatim}
结果寄存器的保存位宽为32位，高96位由输出组合逻辑补零。

\subsection{译码、配置依赖与写回}
精确编码匹配后，译码器将vdot归入向量运算，生成单个微操作并建立两个向量源依赖。实际VL通过重命名后的依赖传递，SEW、LMUL和其他向量控制随指令保存。这样，多条向量配置与运算交错时，每条点积仍使用与自己对应的配置。

旧目的值不参与累加，译码去除相应的旧目的依赖；向量和v0的目的分配、异常信息及最终写回沿用现有处理器机制。选择两槽构建时显式启用\texttt{VDOT\_O2\_ELASTIC=1}和\texttt{VDOT\_T14\_FUSED\_INPUT=1}，并采用对应配置清单。工程同时保留通用单请求路径，便于按构建版本比较和回退。

\subsection{共享VFEX4的派发与门控}
VDot与整数DIV、浮点DIV/SQRT共用VFEX4父级。父级根据选择、实际占用和被选功能单元的ready进行原子派发，使父级和子单元接受的是同一笔请求。VDot两槽排空后可以切换到DIV；DIV在途期间保存自己的ROB身份和占用状态。输出按唯一有效完成来源选择共享写回。

\reportfigure{S2-03}{VFEX4中的派发、资源互斥与共享写回}{10.0cm}{VDot、整数DIV和浮点DIV/SQRT共享父级；各自的真实占用决定派发和时钟保持。}{两槽配置结构示意，依据父级派发与写回实现。}

时钟使能覆盖存活请求、当前输入接受和必要的取消、复位清理。对应验证包含满槽背压、部分取消、ROB环回、关钟后的复位恢复，以及VDot和两类DIV之间的切换。受测的DIV取消修复在匹配请求撤销后释放时钟保持，使父级能够进入空闲并接受后续工作。

\reportfigure{S2-08}{共享资源切换、取消与时钟保持}{8.0cm}{按存活请求保持时钟，在取消或完成后释放相应占用，随后恢复派发。}{门控机制示意，依据真实父级的取消、空闲和恢复检查。}

\section{基础功能与分层验证}
验证首先用独立数学计算确定完整128位结果，再依次检查算术RTL、带握手的功能单元、真实父级和全核执行。算术测试关注有符号乘法、极值和逐级位宽；FU与父级测试加入背压、取消、复位、目的身份和DIV交互；整核协同验证进一步检查译码、架构状态及提交结果。

\reportfigure{S2-04}{从数学结果到整核执行的验证层次}{10.0cm}{模型、算术RTL、FU、真实VFEX4父级和整核各检查对应机制；完整128位结果贯穿各层。}{验证流程示意，依据已有测试入口和结果记录。}

以下三类测试对应赛题2.3.1至2.3.3。图中的输入与响应取自固定八项两槽版本的全核观察记录，按seq、ROB和目的物理寄存器配对；第三阶段说明该版本的架构差分验证。
\clearpage
\subsection{题目2.3.1：正数点积测试}
两组有效输入均为[1,2,3,4,5,6,7,8]，八项平方和为204。历史全核functional-2记录中，seq=1、ROB=0:39、pdest=103在339296周期接受、339298周期响应；实际完整128位结果为\seqsplit{0x000000000000000000000000000000cc}，与预期相同。源高64位为A5哨兵，不参与点积。下图给出对应输入与实际响应摘录。\par
\reportfigure{S2-TEST-01}{正数点积测试：预期与实际结果均为204}{5.0cm}{两个源的有效元素均为1至8；expected与actual完整128位一致，末尾为CC。}{用户保存的原始PNG；functional-2观察器日志字段摘录，保留事务身份、原始行号与SHA256；不是波形截图。}
\subsection{题目2.3.2：负数点积测试}
有效输入为八个−1与八个1，预期总和−8。functional-0的seq=4、ROB=0:12、pdest=143在420291周期接受、420293周期响应；实际完整结果为\seqsplit{0x000000000000000000000000fffffff8}。低32位保存负数补码，高96位为零；该检查不同于第一阶段按字节相加的vadd补码自检。\par
\reportfigure{S2-TEST-02}{负数点积测试：低32位为−8的补码，高96位清零}{5.0cm}{八个−1与八个1点积；实际低32位为FFFFFFF8，高96位为零。}{用户保存的原始PNG；functional-0观察器日志字段摘录，保留事务身份、原始行号与SHA256；不是波形截图。}
\clearpage
\subsection{题目2.3.3：边界值点积测试}
最大正结果使用八个−128与八个−128，每项乘积16384，总和131072。functional-0的seq=5、ROB=0:43、pdest=66在448417周期接受、448419周期响应；实际完整结果为\seqsplit{0x00000000000000000000000000020000}。该例检验乘法和归约在正端点不被截断。\par
\reportfigure{S2-TEST-03}{边界值点积测试：最大正结果131072}{5.0cm}{八对−128相乘求和；实际低32位为00020000，高96位为零。}{用户保存的原始PNG；functional-0观察器日志字段摘录，保留事务身份、原始行号与SHA256；本图只证明所示最大正结果案例。}
零输入用八个0与八个127，实际完整结果为零；functional-0的seq=0、ROB=0:8、pdest=124在311084周期接受、311086周期响应。最小结果的数学预期为8×(−128)×127=−130048，低32位为\seqsplit{0xfffe0400}；本稿未单独展开该项全核输入与响应配对，因此该值只列为预期边界，不新增通过声明。\par
\noindent\textbf{证据来源与判读：}固定0f历史运行，原始记录根目录为\seqsplit{/home/riscv/riscv-lab/software-import-20261026/clean-02/fullcore-run-01/runs/}；正数在\seqsplit{functional-2/run.log}第7030—7034行，负数、最大正结果、零输入分别在\seqsplit{functional-0/run.log}第7927—7931、8400—8404、5814—5818行。按seq、ROB与pdest配对，比较完整128位响应。上述三图是日志字段摘录，不是GTKWave截图；执行响应不等同于架构提交，完整运行终态与差分结论按原归档核对。\par

\subsection{算术数据通路测试汇总}
除前述具体用例外，历史组合算术RTL测试统计为22项定向测试、10000项随机测试，合计10022项全部通过。该汇总与前述全核用例摘录属于不同验证层级，各自保留原测试范围。\par
\reportfigure{S2-TEST-SUMMARY}{VDot算术数据通路定向与随机测试结果}{4.0cm}{定向测试22项、随机测试10000项，合计10022项全部通过。本结果属于算术数据通路验证；全核协同仿真、流水线反压及取消行为由相应测试另行验证。}{用户提供的原始统计截图，完整保留三行PASS输出。本次仅补图，未重新执行测试；后续S2-09综合图中的算术统计为同一组历史结果，不重复计数。}


\reportfigure{S2-09}{正数、负数与边界测试的输入和结果}{12.0cm}{集中对照正数204、负数$-8$和最大正结果131072，观察低32位值及高96位零；算术统计与前述汇总为同一组数据。}{对应测试记录与结果示意，详细响应摘录见本节三张原始观察图。}

在基础数值之外，测试还检查源高半区哨兵、旧目的哨兵、寄存器别名、v0目的、非法配置和向量状态恢复。FU与真实父级通过响应身份、写使能、取消和高位错误注入检查判定链，并在正常运行中核对完整响应与排空。

\section{设计取舍与扩展}
\subsection{两槽结构与结果保存}
两槽结构使乘法和归约能够与相邻请求重叠，面向连续独立点积提高供给能力。将R中的结果保存宽度收窄至32位，利用了高96位固定为零这一架构性质，减少了数据寄存器成本。每笔请求的ROB、目的和控制信息则完整保留，以支持乱序执行、背压和精确取消。

共享父级的优化同时考虑正常完成与取消后的资源释放。方案选择综合功能、周期、标准单元面积和时序路径：一拍算术候选在成本筛选中暴露的父级代价，支持继续使用已有两槽实现作为对应性能表的受测对象。面积和各类路径的具体数据集中列在第三阶段。

\subsection{十六项INT8点积的功能进展}
在八项实现的基础上，团队增加了独立十六项点积候选。该候选以补丁和配套验证快照交付，采用\texttt{MATCH=0x0e00000b}、\texttt{MASK=0xfe00707f}，在e8/m1、VL=16时消费两个128位源的全部16个字节。八项指令继续使用原编码和VL=8语义。

候选复用两组八项乘积及归约模块。P保存16个16位乘积，两组19位部分和相加得到20位精确和，再扩展为低32位结果。十六项和的范围为$[-260096,262144]$，R仍保存32位，写回高96位为零。接收八项指令时，高八个乘积写零，使后级统一完成归约。两槽握手和取消机制沿用原结构。

2026年11月6日，该候选完成了编译开关关闭和开启两种配置下的译码、FU和真实父级验证。启用十六项扩展后，使用新EMU与匹配NEMU完成六例整核差分：2例八项指令、4例十六项指令全部通过。六例覆盖源高半区贡献、有符号正负端点、选定源目的别名、完整128位写回和恢复八项指令。相应FU与父级连续请求场景还核对了2拍首响应与1拍发起间隔。

该扩展的当前成果是指令与执行通路的功能验证。后续将通过完整Q288调用、与较快RVV实现的性能对照，以及同条件下的综合与时序分析，评价其应用价值。本报告的主要周期和成本表对应表中命名的八项受测版本。

\subsection{元素位宽的参数化验证}
另一组独立模块实验把输入元素扩展为INT16，计算8项点积，并使用低64位结果字段、高64位清零的输出布局。该配置的最大正和为$2^{33}$，精确累加至少需要35位。参数化乘法与归约模块完成94项定向输入和10000项随机输入的完整128位结果检查，并验证单请求包装中的反压、取消与复位行为。这组结果展示了从元素位宽到中间精度、结果字段和测试向量的一致扩展方法。

\section{实现代码与测试记录}
以下截图展示乘法与归约实现、主机模型测试，以及两槽控制的关键代码和事务日志。
\evidencefigure{SUP-S2-ARITH}{VDot乘法与归约代码}{15cm}{VDotProducts生成8路有符号乘积，VDotReduce完成归约求和。}
\evidencefigure{SUP-S2-MODEL}{主机参考模型测试结果}{14cm}{主机模型及测试数据检查共35个测试方法通过。}
\evidencefigure{SUP-S2-OVERLAP}{两槽重叠事务日志}{15cm}{clean-02 dense-batch128日志中，10986周期同时发生ACCEPT seq2、MOVE seq1与RESPONSE seq0，分别对应新请求接收、旧P迁移和旧R响应。}
\evidencefigure{SUP-S2-STATE}{两槽寄存器与控制代码}{16cm}{依次展示寄存状态、握手与取消条件，以及读取旧P数据的状态更新。}


\section{源码与证据入口}
\begin{center}
\small
\begin{tabular}{p{0.25\linewidth}p{0.64\linewidth}}
\toprule
内容 & 阅读入口\\
\midrule
指令语义与编码 & \href{https://github.com/cc77-builder/riscv/blob/a0f13456e0c3180fe9c4023f5cbe7874eb5d35a9/team/spec/vdot.md}{指令规格与定向向量}\\
算术及两槽实现 & \href{https://github.com/cc77-builder/riscv/blob/0f12df675378556670afd0dd2153f47fcb330610/src/main/scala/xiangshan/backend/fu/wrapper/VDotElastic.scala}{八项弹性FU源码}及\href{https://github.com/cc77-builder/riscv/blob/0f12df675378556670afd0dd2153f47fcb330610/src/main/scala/xiangshan/backend/fu/vector/VDotDatapath.scala}{乘积与归约源码}\\
构建与版本选择 & \href{https://github.com/cc77-builder/riscv/blob/f5f67e8640d8fda412e8a8946ffb8168f67d79fa/team/docs/development-route-current.md}{工程版本入口与可选配置}\\
周期、面积与STA & \href{https://github.com/cc77-builder/riscv/blob/a0f13456e0c3180fe9c4023f5cbe7874eb5d35a9/team/optimizations/o2/r3-div-gate/PERFORMANCE_COST_RESULT.md}{八项两槽受测版本的数据与原始报告}\\
十六项扩展 & \href{https://github.com/cc77-builder/riscv/blob/f5f67e8640d8fda412e8a8946ffb8168f67d79fa/team/optimizations/o2/dot16/rtl-integration/README.md}{候选结构、补丁与分层验证}、\href{https://github.com/cc77-builder/riscv/blob/f5f67e8640d8fda412e8a8946ffb8168f67d79fa/team/optimizations/o2/dot16/rtl-integration/fullcore/RESULT.md}{六例整核差分结果}\\
元素位宽参数化 & \href{https://github.com/cc77-builder/riscv/blob/a0f13456e0c3180fe9c4023f5cbe7874eb5d35a9/team/docs/t15-parameter-extension.md}{八项INT16独立模块及完整结果验证}\\
\bottomrule
\end{tabular}
\end{center}
工程索引保留完整源码版本、运行输入、原始日志和失败记录。正文以设计机制和结果解释为主，读者可按需要沿上述入口复查实现细节与测量依据。
